%
\subsection*{Problema regular y parámetro espectral}
\label{CH09-B01}
Usamos la forma
\[
 -(p y')'+q y=\lambda w y\quad (a<x<b),
\]
con intervalo finito, $p$ de clase $C^1$, $q,w$ continuas y reales,
y $p,w>0$ en $[a,b]$. Se añaden condiciones homogéneas separadas
de valor, derivada o Robin; en cada extremo,
$\alpha y+\beta y'=0$, con constantes reales no simultáneamente nulas
e independientes de $\lambda$. Un autovalor admite una autofunción
\emph{no idénticamente nula}.

Al resolver, sepárense los casos $\lambda<0$, $\lambda=0$ y
$\lambda>0$ cuando den formas distintas de la EDO. La función constante
no nula puede ser un modo admisible; no debe confundirse con la función
nula. Tampoco cualquier condición de Robin garantiza autovalores
positivos. Recuérdense EDO lineales e integración por partes.
Consulta: solución de \exref{M2-E0054}.

%
\subsection*{Peso, normalización y expansión}
\label{CH09-B02}
El producto escalar correspondiente al problema es
$\langle f,g\rangle_w=\int_a^b wfg\,dx$.
Las autofunciones de autovalores distintos son ortogonales con este peso.
Con $N_n^2=\int_a^b w\phi_n^2\,dx$, la función normalizada es
$e_n=\phi_n/N_n$, y
\[
 f\sim\sum_n c_n\phi_n,\qquad
 c_n=\frac{\int_a^b wf\phi_n\,dx}{N_n^2}.
\]
En la base normalizada, el coeficiente es
$\langle f,e_n\rangle_w=c_nN_n$. Cambiar la normalización cambia los
coeficientes, pero no cada producto coeficiente por autofunción.

Para los problemas regulares y las bases completas aquí utilizadas,
la expansión converge en norma cuadrática ponderada para datos de
cuadrado integrable. Las afirmaciones puntuales necesitan regularidad
adicional. Una representación interior no impone automáticamente los
valores ni las derivadas del dato en los extremos; no se derive la serie
sin justificarlo. Consulta: solución de \exref{M2-E0297}.

%
\subsection*{Contornos físicos y cambio de medida}
\label{CH09-B03}
Si el flujo en la dirección positiva es $J=-k u_x$, con $k>0$, el
flujo saliente del intervalo es $-J(a)=k u_x(a)$ a la izquierda y
$J(b)=-k u_x(b)$ a la derecha. Esos signos deben introducirse antes de
traducir una ley de intercambio a una condición Robin.
Identifíquense $p,q,w$ en la ecuación escrita y compruébense las unidades
de $\lambda$; el peso no es un añadido de notación.

En un cambio monótono y regular $x=\chi(\xi)$, la medida también cambia:
\[
 \int w(x)f(x)g(x)\,dx
 =\int w(\chi(\xi))f(\chi(\xi))g(\chi(\xi))
        |\chi'(\xi)|\,d\xi,
\]
con los límites ordenados de menor a mayor en ambos intervalos.
Transformar las funciones sin el jacobiano puede destruir la
ortogonalidad. Revísense regularidad y positividad en todo el intervalo
cerrado, incluidos los extremos. Consultas: soluciones de
\exref{M2-E0350} y \exref{M2-E0351}.

%
\subsection*{Forzamiento, compatibilidad y límites}
\label{CH09-B04}
Con los mismos contornos homogéneos, escribamos
$-(py')'+qy-\lambda wy=wf$. Proyectar sobre un modo $\phi_n$ da
\[
 (\lambda_n-\lambda)c_n
 =\frac{\langle f,\phi_n\rangle_w}{\|\phi_n\|_w^2}.
\]
Si $\lambda=\lambda_n$, el miembro derecho debe anularse. En los
problemas regulares elementales tratados, con forzamiento regular, la
compatibilidad con el modo resonante permite resolver y deja libre su
coeficiente homogéneo. La resonancia puede dar inexistencia o no
unicidad: compruébese cuál ocurre, antes de calcular una particular.

Si los autovalores se obtienen como raíces, indíquese el intervalo y el
índice de cada raíz. Una aproximación relativa
$\lambda_n/\widetilde\lambda_n\to1$ no obliga al error absoluto a
tender a cero. En límites de familias, úsense cotas con el dominio y los
parámetros especificados; el problema límite puede conservar libertad
sin seleccionar una única solución. Consulta: solución de
\exref{M2-E0088}.
